\secciongrande{Integración del código con las etiquetas html: Integración online o en archivo externo}

\subtituloA{Métodos de integración de javascript}

\definicion{Integración online}{Consiste en incluir el código javascript directamente dentro de las etiquetas %%VERBATIM\_BLOCK\_0%% utilizando el atributo type="texto/javascript"}
\definicion{Integración en archivo externo}{Consiste en referenciar un archivo javascript externo (con extensión .js) mediante el atributo src de la etiqueta %%VERBATIM\_BLOCK\_1%%}
\lista{
  \textbf{Integración online:} %%VERBATIM\_BLOCK\_2%% para evitar problemas de análisis sintáctico.
}
\definicion{Licencias de código}{Permiten añadir licencias, como Creative Commons, para indicar cómo otros desarrolladores pueden utilizar el código.}
\nota{Para el examen}{
¿Cuál es la forma más recomendable de integrar código javascript y por qué?

\textbf{Respuesta:} El archivo externo, debido a la modularidad, facilidad de depuración y rapidez de carga.
}
